\section{Strongly coupled and dusty plasmas}
\label{sec:strong}

The first application profile of the project is a dusty plasma: micron-size
grains, each carrying thousands of elementary charges, suspended in a weakly
ionized gas. The grains are few and individually tracked. Their coupling
parameter $\Gamma$ is large, so the weak-coupling closure of
Section~\ref{sec:plasma} does not apply to them. Three closures are
stacked.

\subsection{First closure: the effective force}

The full system contains electrons, ions, neutrals and fields. The
retained state is the grain positions and velocities. Eliminating the rest
leaves a force on each grain that the grains alone must determine.
\begin{itemize}
  \item The classical closure is an \emph{instantaneous pair law}: a
    screened Coulomb (Yukawa) potential
    $\phi(r)=Q^2e^{-r/\lambda_D}/(4\pi\varepsilon_0r)$, plus confinement
    and gas drag. The state of the grain system is then described by
    $\Gamma$ and the screening parameter $\kappa=a/\lambda_D$.
  \item What it discards: ion flow past a grain produces a wake, so the
    force between two grains is not reciprocal and not a gradient of a
    potential; grain charge depends on the local plasma and fluctuates; the
    wake takes time to form, which is memory; and the force on a grain
    depends on more than one neighbour.
\end{itemize}
In the terms of Section~\ref{sec:kinds} the candidates are an algebraic
pair closure, a many-body (graph) closure, and a memory closure, and the
project's decisive experiment nests them in that order. The eliminated
plasma has its own dynamics and is not in equilibrium, so the
fluctuation--dissipation relation \eqref{eq:gle} need not hold for the
effective grain dynamics. The data-driven force inference of the Burton
laboratory \cite{Yu25} is an instance of learning this closure.

\subsection{Second closure: correlations}

Given a pair law, the equilibrium hierarchy relates the pair distribution
$g(r)$ to the triplet distribution. Liquid-state theory calls its
approximations \emph{closures}, in exactly the sense used here
\cite{HansenMcDonald}.
\begin{itemize}
  \item \textbf{Kirkwood superposition:}
    $g_3(1,2,3)\approx g(1,2)\,g(2,3)\,g(1,3)$, the factorization of the
    triplet function into pairs.
  \item \textbf{Integral-equation closures.} The Ornstein--Zernike relation
    $h=c+n\,c*h$ defines the direct correlation function $c$ from
    $h=g-1$. It is an identity with two unknowns. A second relation closes
    it:
    \[
      \text{hypernetted chain:}\quad g=\exp(-\phi/T+h-c),
      \qquad
      \text{Percus--Yevick:}\quad c=(1-e^{\phi/T})\,g .
    \]
    The first is accurate for long-range, soft potentials such as Coulomb
    and Yukawa; the second for short-range, hard ones.
\end{itemize}
For the dynamics, the BBGKY closure $g_2\to0$ fails outright at
$\Gamma\gtrsim1$: the system is a liquid, and for the one-component plasma
it freezes at $\Gamma\approx175$. Kinetic descriptions that keep
correlations use them as input: local-field corrections to the mean field,
and the quasi-localized charge approximation, in which particles oscillate
in a slowly changing cage described by $g(r)$ \cite{GoldenKalman00}.

\subsection{Third closure: moments with memory}

At the fluid level a strongly coupled liquid supports shear waves at short
wavelength, which a Newtonian stress cannot. The standard closure is
viscoelastic: the stress relaxes over a time $\tau_m$,
\begin{equation}
  \Bigl(1+\tau_m\frac{\partial}{\partial t}\Bigr)\boldsymbol\pi
  =-\mu\bigl(\nabla\bu+\nabla\bu^{\mathsf T}
     -\tfrac23(\nabla\cdot\bu)\mathbf I\bigr),
\end{equation}
which is Cattaneo's law (Example~\ref{ex:cattaneo}) for momentum, a
one-pole memory kernel. It gives viscous behaviour for
$\omega\tau_m\ll1$ and elastic behaviour for $\omega\tau_m\gg1$. This
generalized-hydrodynamic model was applied to dusty plasmas by Kaw and Sen
\cite{KawSen98}. Here the memory is not a refinement: without it the
closure has the wrong qualitative behaviour.

\subsection{What transfers}

The dusty-plasma profile is where the discarded information is
\emph{observable}: grain trajectories give two- and three-particle
statistics directly, so the factorization defect $g_2$, and the size of the
memory and many-body parts of the force, can be measured instead of
assumed. That is the reason it is paired with the canonical tracks, in
which the same quantities are controlled by theorems.
